% Ngân hàng tự sinh (Python)
\begin{bt}%[0D0N1-1]
% nguon: Python sinh từ câu mẫu 5199922e
% mau: 5199922e
	Gieo ngẫu nhiên cùng lúc  $6$ đồng xu phân biệt cân đối và đồng chất. Số phần tử của không gian mẫu là gì?
	\loigiai{
		Trong một lần gieo, mỗi đồng  xu có $6$ khả năng xuất hiện.\\
		Vậy số phần tử của không gian mẫu trong phép thử trên là $n(\Omega)=6\cdot 6=36$.
	}
\end{bt}

\begin{bt}%[0D0H1-1]
% nguon: Python sinh (số lớn hơn) từ câu mẫu 5199922e
% mau: 5199922e nang_cao
	Gieo ngẫu nhiên cùng lúc  $4$ đồng xu phân biệt cân đối và đồng chất. Số phần tử của không gian mẫu là gì?
	\loigiai{
		Trong một lần gieo, mỗi đồng  xu có $4$ khả năng xuất hiện.\\
		Vậy số phần tử của không gian mẫu trong phép thử trên là $n(\Omega)=4\cdot 4=16$.
	}
\end{bt}

\begin{bt}%[0D0H1-3]
% nguon: Python sinh từ câu mẫu 811cd64d
% mau: 811cd64d
	Gieo một đồng tiền liên tiếp $4$ lần. Số phần tử của không gian mẫu $n(\Omega)$ là gì?
	\loigiai{
		Số phần  tử của không gian mẫu $n(\Omega)=4\cdot 4=16$.
	}
\end{bt}

\begin{bt}%[0D0V1-3]
% nguon: Python sinh (số lớn hơn) từ câu mẫu 811cd64d
% mau: 811cd64d nang_cao
	Gieo một đồng tiền liên tiếp $6$ lần. Số phần tử của không gian mẫu $n(\Omega)$ là gì?
	\loigiai{
		Số phần  tử của không gian mẫu $n(\Omega)=6\cdot 6=36$.
	}
\end{bt}

\begin{bt}%[0D0H2-5]
% nguon: Python sinh từ câu mẫu 2ca6f646
% mau: 2ca6f646
	Gieo  một con xúc xắc hai lần liên tiếp. Số khả năng thuận lợi để \lq\lq \,Tích số chấm ở hai lần gieo là một số nhỏ hơn $8$\,\rq\rq \ là bao nhiêu?
	\loigiai{
		Số phần tử của không gian mẫu là $n\left(\Omega\right)=8\cdot 8=64$.\\
		Gọi $A$ là biến cố:  \lq\lq \,Tích số chấm ở hai lần gieo là một số nhỏ hơn $8$\,\rq\rq.\\
		Ta có $A=\left\{(1;1); (1;2); (2;1); (1;3); (3;1); (1;4); (4;1); (1;5);(5;1); (2;2)\right\}$.\\
		Số kết  quả thuận lợi cho biến cố $A$ là $n(A)=10$.\\
	}
\end{bt}

\begin{bt}%[0D0V2-5]
% nguon: Python sinh (số lớn hơn) từ câu mẫu 2ca6f646
% mau: 2ca6f646 nang_cao
	Gieo  một con xúc xắc hai lần liên tiếp. Số khả năng thuận lợi để \lq\lq \,Tích số chấm ở hai lần gieo là một số nhỏ hơn $18$\,\rq\rq \ là bao nhiêu?
	\loigiai{
		Số phần tử của không gian mẫu là $n\left(\Omega\right)=18\cdot 18=324$.\\
		Gọi $A$ là biến cố:  \lq\lq \,Tích số chấm ở hai lần gieo là một số nhỏ hơn $18$\,\rq\rq.\\
		Ta có $A=\left\{(1;1); (1;2); (2;1); (1;3); (3;1); (1;4); (4;1); (1;5);(5;1); (2;2)\right\}$.\\
		Số kết  quả thuận lợi cho biến cố $A$ là $n(A)=10$.\\
	}
\end{bt}
